% =============================================================
%  CHAPTER 2 - ANALYSIS & PROJECT PLANNING
% =============================================================

\chapteroverview{Analysis and Project Planning}{
  \item Introduction
  \item Identification of the Actors
  \item Functional Requirements
  \item Non-Functional Requirements
  \item Global Use Case Diagram
  \item Product Backlog
  \item Releases and Sprints
  \item Provisional Schedule (Gantt)
  \item Conclusion
}

\section{Introduction}

After framing the project, this chapter focuses on understanding \textit{what} the ADVANCIA platform must do and \textit{how} the work is organized. We start by identifying every actor of the system, we then formalize the functional and non-functional requirements, we draw the global use case diagram and finally we transform the whole into a structured plan of releases, sprints and a Gantt schedule.

% ---------- 2.2 Actors ----------
\section{Identification of the Actors}

The ADVANCIA platform is a multi-actor system. Four main actors interact with the system, each one with a clearly defined scope of responsibilities.

\begin{table}[H]
\centering
\renewcommand{\arraystretch}{1.35}
\begin{tabularx}{\textwidth}{|>{\columncolor{softblue}}l|X|}
\hline
\rowcolor{primary}
\textbf{\color{white} Actor} & \textbf{\color{white} Role and responsibilities}\\
\hline
\textbf{User} &
Creates and manages his account, browses and reserves trainings, manages his profile (photo, avatar, password) and consults his personal dashboard. \\\hline
\textbf{Trainer} &
Handles the trainings he is responsible for, follows enrolled users, animates sessions and updates pedagogical content. \\\hline
\textbf{Administrator} &
Manages users, validates or refuses reservations, imports and exports data (Excel, PDF), exploits a smart dashboard and an integrated chatbot. \\\hline
\textbf{Super-Administrator} &
Has the highest level of authority : manages users, admins, trainers, trainings and courses, sends notifications and supervises the whole platform. \\\hline
\end{tabularx}
\caption{Actors of the ADVANCIA platform}
\label{tab:actors}
\end{table}

% ---------- 2.3 Functional ----------
\section{Functional Requirements}

The functional requirements describe \textit{what} the system must offer. They are grouped by actor.

\subsection{User}
\begin{itemize}[leftmargin=2em,itemsep=0.2em]
  \item Create an account and authenticate.
  \item Manage the personal profile : photo, avatar, password, contact info.
  \item Browse and search the catalog of trainings.
  \item Reserve a training session.
  \item Consult a simple personal dashboard with the essential indicators.
  \item Receive notifications about reservations and account events.
\end{itemize}

\subsection{Trainer}
\begin{itemize}[leftmargin=2em,itemsep=0.2em]
  \item Authenticate and manage profile.
  \item Consult the list of trainings and sessions assigned to him.
  \item Track enrolled users.
  \item Update pedagogical resources of his courses.
\end{itemize}

\subsection{Administrator}
\begin{itemize}[leftmargin=2em,itemsep=0.2em]
  \item Manage users (create, update, suspend, delete).
  \item Import data (e.g. CSV/Excel of users) and export data in Excel or PDF.
  \item Use the smart dashboard with KPIs and charts.
  \item Use the integrated chatbot for assistance and quick queries.
  \item Manage personal profile.
\end{itemize}

\subsection{Super-Administrator}
\begin{itemize}[leftmargin=2em,itemsep=0.2em]
  \item All capabilities of the Administrator.
  \item Manage admins and trainers.
  \item Manage trainings and courses (CRUD).
  \item Accept or refuse reservations.
  \item Send notifications to users or administrators.
  \item Change password and supervise security policies.
\end{itemize}

% ---------- 2.4 Non functional ----------
\section{Non-Functional Requirements}

\begin{table}[H]
\centering
\renewcommand{\arraystretch}{1.35}
\begin{tabularx}{\textwidth}{|>{\columncolor{softblue}}l|X|}
\hline
\rowcolor{primary}
\textbf{\color{white} Requirement} & \textbf{\color{white} Description}\\
\hline
Security & Authentication via JWT, password hashing, role-based access control, HTTPS. \\\hline
Performance & Response time below 2 seconds for the main pages, smooth navigation. \\\hline
Usability & Modern, clear and responsive UI accessible from desktop, tablet and mobile. \\\hline
Reliability & Daily database backup and graceful error handling. \\\hline
Scalability & Modular architecture allowing the addition of new modules without disruption. \\\hline
Maintainability & Clean code, documented modules, version-controlled with Git. \\\hline
Portability & Cross-browser compatibility (Chrome, Firefox, Edge). \\\hline
\end{tabularx}
\caption{Non-functional requirements}
\label{tab:nonfunctional}
\end{table}

% ---------- 2.5 Use case diagram ----------
\section{Global Use Case Diagram}

The following diagram synthesizes the main use cases of the ADVANCIA platform and the relationships between actors and use cases.

\begin{figure}[H]
\centering
\resizebox{\textwidth}{!}{%
\begin{tikzpicture}[
  every node/.style={font=\small},
  uc/.style={ellipse, draw=primary, fill=softblue, minimum width=3cm,
             minimum height=0.7cm, align=center, inner sep=2pt},
  actor/.style={align=center, font=\small\bfseries}
]
  % System boundary
  \draw[draw=primary, thick, rounded corners=8pt] (-1,-7.5) rectangle (10,5);
  \node[font=\bfseries\color{primary}] at (4.5,4.6) {ADVANCIA System};

  % Use cases
  \node[uc] (auth)   at (4.5, 4)  {Authenticate};
  \node[uc] (regis)  at (1.0, 3)  {Register};
  \node[uc] (profile)at (4.5, 3)  {Manage profile};
  \node[uc] (browse) at (1.0, 2)  {Browse trainings};
  \node[uc] (reserve)at (4.5, 2)  {Reserve training};
  \node[uc] (dash)   at (8.0, 3)  {View dashboard};
  \node[uc] (course) at (8.0, 2)  {Manage courses};
  \node[uc] (mu)     at (1.0, 0.5){Manage users};
  \node[uc] (ma)     at (4.5, 0.5){Manage admins};
  \node[uc] (mt)     at (8.0, 0.5){Manage trainings};
  \node[uc] (accept) at (1.0,-1)  {Accept / refuse\\reservation};
  \node[uc] (notif)  at (4.5,-1)  {Send notifications};
  \node[uc] (chat)   at (8.0,-1)  {Use chatbot};
  \node[uc] (imp)    at (1.0,-2.5){Import data};
  \node[uc] (exp)    at (4.5,-2.5){Export Excel/PDF};
  \node[uc] (track)  at (8.0,-2.5){Track learners};

  % Actors (left: User, Trainer)
  \node[actor] (user) at (-3,3)
       {\tikz\draw[primary, thick] (0,0) circle (0.18) (0,-0.18) -- (0,-0.7)
         (-0.4,-0.4) -- (0.4,-0.4) (-0.25,-1.05) -- (0,-0.7) -- (0.25,-1.05);\\User};
  \node[actor] (trainer) at (-3,0)
       {\tikz\draw[primary, thick] (0,0) circle (0.18) (0,-0.18) -- (0,-0.7)
         (-0.4,-0.4) -- (0.4,-0.4) (-0.25,-1.05) -- (0,-0.7) -- (0.25,-1.05);\\Trainer};

  % Actors (right: Admin, Super-Admin)
  \node[actor] (admin) at (12,3)
       {\tikz\draw[primary, thick] (0,0) circle (0.18) (0,-0.18) -- (0,-0.7)
         (-0.4,-0.4) -- (0.4,-0.4) (-0.25,-1.05) -- (0,-0.7) -- (0.25,-1.05);\\Admin};
  \node[actor] (super) at (12,-1)
       {\tikz\draw[primary, thick] (0,0) circle (0.18) (0,-0.18) -- (0,-0.7)
         (-0.4,-0.4) -- (0.4,-0.4) (-0.25,-1.05) -- (0,-0.7) -- (0.25,-1.05);\\Super-Admin};

  % Links - User
  \draw (user) -- (auth);
  \draw (user) -- (regis);
  \draw (user) -- (profile);
  \draw (user) -- (browse);
  \draw (user) -- (reserve);
  \draw (user) -- (dash);

  % Links - Trainer
  \draw (trainer) -- (auth);
  \draw (trainer) -- (course);
  \draw (trainer) -- (track);

  % Links - Admin
  \draw (admin) -- (auth);
  \draw (admin) -- (mu);
  \draw (admin) -- (dash);
  \draw (admin) -- (chat);
  \draw (admin) -- (imp);
  \draw (admin) -- (exp);

  % Links - Super-Admin
  \draw (super) -- (mu);
  \draw (super) -- (ma);
  \draw (super) -- (mt);
  \draw (super) -- (accept);
  \draw (super) -- (notif);
  \draw (super) -- (chat);
\end{tikzpicture}}
\caption{Global use case diagram of the ADVANCIA platform}
\label{fig:globaluc}
\end{figure}

% ---------- 2.6 Product Backlog ----------
\section{Product Backlog}

The product backlog gathers, in priority order, every user story to be developed. The story points (SP) follow a Fibonacci scale.

\begin{longtable}{|>{\columncolor{softblue}}p{0.9cm}|p{8.5cm}|p{2.3cm}|p{1.1cm}|}
\hline
\rowcolor{primary}
\textbf{\color{white} ID} & \textbf{\color{white} User Story} & \textbf{\color{white} Actor} & \textbf{\color{white} SP}\\
\hline
\endfirsthead
\hline
\rowcolor{primary}
\textbf{\color{white} ID} & \textbf{\color{white} User Story} & \textbf{\color{white} Actor} & \textbf{\color{white} SP}\\
\hline
\endhead
US01 & As a user, I want to create an account so I can access the platform. & User & 3 \\\hline
US02 & As a user, I want to authenticate securely. & User & 3 \\\hline
US03 & As a user, I want to manage my profile (photo, avatar, password). & User & 5 \\\hline
US04 & As a user, I want a personal dashboard. & User & 5 \\\hline
US05 & As a user, I want to browse and search trainings. & User & 5 \\\hline
US06 & As a user, I want to reserve a training session. & User & 5 \\\hline
US07 & As a trainer, I want to consult my assigned trainings. & Trainer & 3 \\\hline
US08 & As a trainer, I want to track my enrolled learners. & Trainer & 5 \\\hline
US09 & As an admin, I want to manage users (CRUD). & Admin & 8 \\\hline
US10 & As an admin, I want to import data from Excel. & Admin & 5 \\\hline
US11 & As an admin, I want to export data to Excel and PDF. & Admin & 5 \\\hline
US12 & As an admin, I want a smart dashboard. & Admin & 8 \\\hline
US13 & As an admin, I want an integrated chatbot. & Admin & 8 \\\hline
US14 & As a super-admin, I want to manage admins and trainers. & Super-Admin & 8 \\\hline
US15 & As a super-admin, I want to manage trainings and courses. & Super-Admin & 8 \\\hline
US16 & As a super-admin, I want to accept or refuse reservations. & Super-Admin & 5 \\\hline
US17 & As a super-admin, I want to send notifications. & Super-Admin & 5 \\\hline
US18 & As any actor, I want a responsive and modern UI. & All & 5 \\\hline
\caption{Product backlog of the ADVANCIA platform}
\label{tab:backlog}
\end{longtable}

% ---------- 2.7 Releases and sprints ----------
\section{Releases and Sprints}

The product backlog has been broken down into \textbf{three releases}, each one delivering a tangible set of value at the end of its lifecycle.

\begin{table}[H]
\centering
\renewcommand{\arraystretch}{1.35}
\begin{tabularx}{\textwidth}{|>{\columncolor{softblue}}l|X|l|}
\hline
\rowcolor{primary}
\textbf{\color{white} Release} & \textbf{\color{white} Scope} & \textbf{\color{white} Sprints}\\
\hline
Release 1 & Authentication \& User management (US01--US04, US18) & S1, S2 \\\hline
Release 2 & Training catalog \& Course management (US05, US07, US15) & S3, S4 \\\hline
Release 3 & Reservation, Notifications, Dashboard \& Chatbot (US06, US08--US14, US16, US17) & S5, S6 \\\hline
\end{tabularx}
\caption{Decomposition into releases and sprints}
\label{tab:releases}
\end{table}

\subsection{Distribution of effort}

\begin{figure}[H]
\centering
\begin{tikzpicture}
\begin{axis}[
  width=12cm, height=5.5cm,
  ybar stacked,
  bar width=22pt,
  ylabel={Story points},
  symbolic x coords={Release 1, Release 2, Release 3},
  xtick=data,
  ymin=0, ymax=70,
  legend style={at={(0.5,-0.20)}, anchor=north, legend columns=-1, draw=none},
  enlarge x limits=0.35,
  nodes near coords style={font=\small\color{white}\bfseries},
  every axis label/.append style={font=\small},
  tick label style={font=\small}
]
  \addplot[fill=primary]    coordinates {(Release 1,16) (Release 2,16) (Release 3,21)};
  \addplot[fill=secondary]  coordinates {(Release 1,5)  (Release 2,8)  (Release 3,13)};
  \addplot[fill=accent]     coordinates {(Release 1,3)  (Release 2,5)  (Release 3,8) };
  \legend{Core features, Advanced features, Polishing}
\end{axis}
\end{tikzpicture}
\caption{Distribution of story points across releases}
\label{fig:effort}
\end{figure}

% ---------- 2.8 Gantt ----------
\section{Provisional Schedule (Gantt Chart)}

\begin{figure}[H]
\centering
\resizebox{\textwidth}{!}{%
\begin{ganttchart}[
  hgrid, vgrid,
  bar/.append style={fill=primary, draw=primary},
  bar incomplete/.append style={fill=softblue, draw=primary},
  group/.append style={fill=secondary, draw=secondary},
  milestone/.append style={fill=accent, draw=accent},
  bar height=0.6,
  group label font=\bfseries\color{white},
  bar label font=\small,
  title label font=\bfseries\color{primary},
  title height=1,
  x unit=0.55cm,
  y unit chart=0.55cm
]{1}{20}
  \gantttitle{ADVANCIA Project -- Provisional Planning}{20} \\
  \gantttitlelist{1,...,20}{1} \\
  \ganttgroup{Analysis \& Design}{1}{4} \\
  \ganttbar{Requirements}{1}{2} \\
  \ganttbar{UML modeling}{2}{4} \\
  \ganttgroup{Release 1}{4}{8} \\
  \ganttbar{Sprint 1 -- Auth}{4}{6} \\
  \ganttbar{Sprint 2 -- User mgmt}{6}{8} \\
  \ganttgroup{Release 2}{8}{13} \\
  \ganttbar{Sprint 3 -- Catalog}{8}{10} \\
  \ganttbar{Sprint 4 -- Courses}{11}{13} \\
  \ganttgroup{Release 3}{13}{18} \\
  \ganttbar{Sprint 5 -- Reservations}{13}{15} \\
  \ganttbar{Sprint 6 -- Dashboard \& Chatbot}{16}{18} \\
  \ganttgroup{Tests \& Documentation}{18}{20} \\
  \ganttbar{Final tests}{18}{19} \\
  \ganttbar{Documentation}{19}{20} \\
  \ganttmilestone{Final delivery}{20}
\end{ganttchart}}
\caption{Gantt chart of the ADVANCIA project}
\label{fig:gantt}
\end{figure}

\section{Conclusion}

This second chapter has clarified the perimeter of the platform : its four actors, its functional and non-functional requirements, the global use cases, and the structuring of the work into three releases driven by a Gantt schedule. The next chapters will detail each release, sprint by sprint, until the final integrated solution.
